
% ════════════════════════════════════════════
% 附录（11.附录.tex）—— 写作提示词见本文件末尾【本章节写作提示词】。
% ════════════════════════════════════════════
\newpage
\makeatletter
\let\@oldaddcontentsline\addcontentsline
\renewcommand{\addcontentsline}[3]{}

\section*{附录}
\let\addcontentsline\@oldaddcontentsline
\makeatother

\renewcommand{\arraystretch}{1.2}
\begin{longtable}{>{\centering\arraybackslash}p{0.30\textwidth}>{\raggedright\arraybackslash}p{0.60\textwidth}}
  \caption{附件说明表}
  \label{tab:appendix} \\
  \toprule
  文件 & 说明 \\
  \midrule
  \endfirsthead
  \caption{附件说明表（续）} \\
  \toprule
  文件 & 说明 \\
  \midrule
  \endhead
  \bottomrule
  \endfoot
  求解/\allowbreak{}问题一/\allowbreak{}问题一.py & 问题一可观测区域与最优观测站求解代码 \\
  求解/\allowbreak{}问题二/\allowbreak{}问题二.py & 问题二交会定位与并发调度求解代码 \\
  求解/\allowbreak{}问题三/\allowbreak{}问题三.py & 问题三误差传播与蒙特卡洛验证代码 \\
  求解/\allowbreak{}问题四/\allowbreak{}问题四.py & 问题四扰动重调度与场景比较代码 \\
  求解/\allowbreak{}\_common.py & 公共绘图与文件工具（字体、高清出图、CSV 读写等） \\
  求解/\allowbreak{}\_geometry.py & 公共几何模块（建筑/车辆/告警参数与可观测性判定） \\
  求解/\allowbreak{}问题一/\allowbreak{}结果/\allowbreak{}问题一\_\allowbreak{}最优观测站参数.csv & 问题一最优站参数结果 \\
  求解/\allowbreak{}问题二/\allowbreak{}结果/\allowbreak{}问题二\_\allowbreak{}1\_\allowbreak{}T2最优站对.csv & 问题二 T2 最优站对结果 \\
  求解/\allowbreak{}问题二/\allowbreak{}结果/\allowbreak{}问题二\_\allowbreak{}2\_\allowbreak{}最优调度时间轴.csv & 问题二并发调度完整时间轴 \\
  求解/\allowbreak{}问题三/\allowbreak{}结果/\allowbreak{}问题三\_\allowbreak{}误差传播结果.csv & 问题三误差传播与参数约束结果 \\
  求解/\allowbreak{}问题四/\allowbreak{}结果/\allowbreak{}问题四\_\allowbreak{}场景结果.csv & 问题四各场景完成率与最晚时间 \\
  求解/\allowbreak{}问题一/\allowbreak{}图片/\allowbreak{}图1-1至图1-22 & 问题一可观测区域与最优站可视化（22 张） \\
  求解/\allowbreak{}问题二/\allowbreak{}图片/\allowbreak{}图2-1至图2-23 & 问题二交会定位与调度可视化（23 张） \\
  求解/\allowbreak{}问题三/\allowbreak{}图片/\allowbreak{}图3-1至图3-22 & 问题三误差传播与灵敏度可视化（22 张） \\
  求解/\allowbreak{}问题四/\allowbreak{}图片/\allowbreak{}图4-1至图4-20 & 问题四重调度与场景对比可视化（20 张） \\
\end{longtable}
